\chapter{Le choix du guide}

\sessioninfo{Session 7}{Mardi 29 septembre 2026}

Devant la demeure de Wakanga O'Tamu, les quatre aventuriers hésitaient encore sur la direction à prendre. La jungle les attendait au-delà des murs de Port Nyanzaru, mais il leur manquait toujours quelqu'un capable de les y conduire. Avant qu'ils ne puissent se mettre en route, le domestique de Wakanga les rattrapa, une lettre à la main.

Son maître avait oublié de la leur remettre. Il s'agissait, semblait-il, du dernier rapport préparé par Iolka et Nera Vesper. Les deux sœurs protégeaient leurs messages afin que seuls leurs véritables destinataires puissent en saisir le sens. Le groupe reçut donc une lettre parfaitement lisible dont le contenu, lui, leur échappait presque entièrement.

\illustration[0.9\textwidth]{2026-09-29/lettre-paire-T22.png}
{\textit{Le dernier rapport attribué à Iolka et Nera Vesper.}}
{lettre_paire_cycle_22}

Le texte évoquait un temple, des ruines à l'ouest, une guide, des tombes et une étrange main qui y fleurissait. Il parlait aussi d'une porte bleue, d'une cloche et d'un fleuve, avant de recommander de ne pas écrire avant le prochain cycle. La signature, réduite à un simple « N. \& I. », renforçait l'idée qu'il provenait bien de la Paire. Sans la clé de lecture, cependant, impossible de distinguer les informations véritables des détours destinés à tromper un lecteur indiscret.

Ils rangèrent soigneusement le message. Rien ne garantissait que Wakanga possédait lui-même la clé nécessaire, ni qu'une nouvelle discussion avec lui leur apporterait une réponse immédiate. Pour l'heure, le choix d'un guide restait le problème le plus urgent.

\scenebreak

Les prospectus consultés à la maison de repos de Kaya présentaient plusieurs candidats. Certains réclamaient le tarif habituel, tandis que d'autres proposaient de réduire ou d'abandonner leurs honoraires en échange d'un service. Kawasha et Kupalué avaient d'abord retenu l'attention du groupe : leur lutte contre les morts-vivants semblait pouvoir rejoindre les préoccupations de l'expédition.

Encore fallait-il les trouver.

Le domestique leur indiqua plusieurs endroits où obtenir des renseignements. Les guides pouvaient être rencontrés chez Jobal, dans les tavernes qu'ils fréquentaient ou près de la capitainerie, là où s'organisaient certaines expéditions. Le groupe choisit cette dernière destination et redescendit vers le port.

Près de l'entrée, les aventuriers retrouvèrent les annonces déjà aperçues chez Kaya. Un prospectus supplémentaire avait toutefois été glissé parmi les autres. De facture plus pauvre que les documents officiels, il vantait les services de River Mist et Flask of Wine tout en insultant ouvertement Jobal et les guides placés sous son autorité.

\illustration[0.8\textwidth]{2026-09-29/098-ph08.png}
{\textit{Le prospectus de River Mist et Flask of Wine découvert à la capitainerie.}}
{prospectus_river_mist_flask_of_wine}

\begin{quote}\itshape
Méfiez-vous de Jobal et de ses incapables de sous-fifres ! Ce prince marchand est aussi digne de confiance qu'un sac de serpents, tout comme la plupart de ses « guides ». Si vous nous choisissez, nous ne vous égarerons jamais.

« River et Flask sont à la hauteur ! » Voilà notre devise ! Chult est notre terrain de jeu. Nous sommes rapides, fiables et savons ce que nous faisons.

Les sous-fifres de Jobal exigent cinq pièces d'or par jour et trente jours payés d'avance. Nous serons heureux de vous servir de guides pour seulement quatre pièces d'or par jour, sans aucun paiement anticipé. Mieux encore, nous renoncerons entièrement à nos honoraires en échange d'une part égale de tous les trésors que nous vous aiderons à découvrir.
\end{quote}

L'offre paraissait avantageuse, mais sa présentation inspira surtout de la méfiance. Le prospectus semblait avoir été ajouté sans autorisation, peut-être après que d'autres annonces eurent été retirées pour lui faire de la place. Les aventuriers préféraient encore traiter avec un guide reconnu. Wakanga leur avait paru honnête et n'aurait eu aucun intérêt à les confier sciemment à quelqu'un d'incapable de mener à bien leurs missions.

Ils poussèrent donc la porte de la capitainerie. L'intérieur était parfaitement tenu, équipé de girouettes et de différents instruments destinés à suivre le vent et les conditions de navigation. Derrière le comptoir se tenait un Sang-dragon doré âgé, vêtu d'une élégante toge blanche. Deux longs appendices descendaient de son museau comme une moustache lorsqu'il les accueillit chaleureusement.

À l'évocation de Kawasha, son expression resta bienveillante. C'était un excellent guide, mais il ne se trouvait pas actuellement à Port Nyanzaru. On pouvait généralement le trouver près du fort Beluarian, à plusieurs jours de marche vers l'est, lorsqu'il ne parcourait pas déjà la jungle avec ses alliés.

Destos déplia alors la carte de Chult pour que leur interlocuteur leur indique la région. Le Sang-dragon admira le document, puis leur recommanda aussitôt de se montrer beaucoup plus prudents. Une carte aussi précise valait une fortune sur une île encore largement inexplorée. Certaines personnes pourraient voler ou même tuer pour s'en emparer.

Le trajet jusqu'au fort paraissait relativement simple en suivant la côte, mais il imposerait plusieurs jours de voyage sans guide. Plus on s'enfonçait vers l'intérieur des terres, plus les risques de croiser des morts-vivants augmentaient. Le groupe n'était pas venu chercher quelqu'un capable de le guider pour commencer par traverser seul une partie de l'île.

Il leur fallait un autre choix.

Le Sang-dragon leur confirma que plusieurs guides se trouvaient actuellement en ville. Azaka Stormfang fréquentait notamment le Lézard Tonitruant, une taverne devant laquelle ils étaient déjà passés. Après avoir remercié leur informateur, les aventuriers reprirent donc les rues désormais un peu plus familières de Port Nyanzaru.

\scenebreak

Le Lézard Tonitruant offrait une ambiance bien plus animée que la maison de repos de Kaya. Au milieu des conversations et des chopes, un homme dédicaçait avec enthousiasme son dernier guide consacré aux monstres devant une file de curieux. Les quatre compagnons laissèrent cette agitation de côté et observèrent plutôt les clients installés au comptoir.

Une femme correspondant au portrait d'Azaka s'y trouvait effectivement. Grande et musclée, la peau sombre et les cheveux bruns, elle portait une large boucle d'oreille et une cicatrice de griffure marquait son menton. Sa posture suffisait à révéler une longue habitude de la marche et du combat. Lorsque les aventuriers l'abordèrent, elle délaissa sa bière et leur demanda simplement où ils comptaient aller.

Le groupe lui expliqua vouloir suivre la rivière vers l'intérieur de Chult, puis bifurquer pour gagner la région où pourrait se trouver le Gardien Bouclier. Il resta volontairement vague sur la nature exacte de l'objet recherché. Azaka ne se montra pas curieuse : leurs affaires leur appartenaient, son rôle serait seulement de les mener jusqu'à leur destination.

Elle affirma avoir grandi dans la jungle et savoir y trouver de la nourriture, choisir un campement sûr, allumer un feu et s'orienter sans difficulté. Le tarif réglementaire était bien de cinq pièces d'or par jour, avec trente jours réglés à l'avance. Une telle somme représentait presque toutes les économies réunies du groupe.

Azaka proposa heureusement une autre solution. À Firefinger se trouvait un masque de bois qu'elle désirait récupérer, mais un ptéranodon maléfique occupait les lieux. Si les aventuriers l'aidaient à reprendre le masque, elle renoncerait entièrement à ses honoraires et paierait elle-même ce qu'elle devait à Jobal.

Firefinger les entraînerait plus à l'est que leur destination initiale. En canoë et avec une météo favorable, le détour demanderait trois ou quatre jours pour descendre la rivière, autant pour revenir, ainsi qu'une journée sur place. L'alternative restait malgré tout préférable à un paiement de cent cinquante pièces d'or. Elle leur permettrait aussi de découvrir la jungle aux côtés d'une guide avant de reprendre leurs deux missions principales.

Les aventuriers demandèrent quelques instants pour discuter à l'extérieur. Les salamandres de feu évoquées par un autre guide leur inspiraient moins confiance qu'un unique ptéranodon, même qualifié de maléfique. Quant à rejoindre Kawasha au fort Beluarian, cela leur ferait perdre plusieurs jours en les obligeant à voyager seuls.

Le choix fut arrêté.

Ils retournèrent auprès d'Azaka et lui donnèrent leur parole. Elle les retrouverait le lendemain matin à la maison de repos de Kaya, à condition qu'ils se procurent auparavant un canoë et le nécessaire pour survivre loin de la ville. Avant le départ, elle devait de son côté régler ses taxes auprès de Jobal.

\scenebreak

Azaka leur dressa une première liste de précautions. Ils auraient besoin de réserves, de protections contre les insectes et de moyens de lutter contre les poisons et les maladies. Surtout, ils ne pourraient pas transporter toute l'eau nécessaire à l'expédition. L'eau de la rivière n'étant pas potable, il leur faudrait recueillir la pluie ou compter sur la magie de Destos.

Le clerc assura pouvoir créer de l'eau en cas de besoin. Ses compagnons lui firent néanmoins remarquer que tout leur approvisionnement reposerait alors sur lui. Un récupérateur de pluie constituerait une sécurité raisonnable si Destos venait à être blessé, malade ou simplement incapable d'utiliser ses pouvoirs.

La jungle ne pardonnait pas l'improvisation. Une simple blessure, une piqûre ou une maladie bénigne pouvait devenir mortelle à plusieurs jours de la ville. Azaka leur parla notamment de la fièvre de la brume bleue, capable de provoquer des hallucinations pendant des heures. Elle la considérait encore comme l'une des maladies les plus sympathiques de Chult.

Cette précision ne rassura personne.

Le groupe quitta la taverne pour rejoindre le Grand Souk. Contrairement au tumulte du bazar rouge, ce marché accueillait une clientèle fortunée dans une atmosphère plus calme. Les marchandises luxueuses côtoyaient les bijoux, les armes étincelantes et le matériel d'exploration vendu sous l'autorité d'Ekene-Afa.

Kael remarqua plusieurs articles inhabituels sur un étal : des blocs odorants à la texture grasse, de grandes toiles repliées autour de tiges de bois et une arme portée à la ceinture par de nombreux Chultiens. Le marchand leur présenta d'abord ses répulsifs contre les insectes. Une lotion plus efficace pouvait protéger celui qui l'appliquait durant une journée entière, même sous la pluie.

Les toiles formaient, une fois déployées, un large entonnoir capable de recueillir rapidement l'eau des averses tropicales. Quant à l'arme, il s'agissait d'une Yklwa, une arme traditionnelle de Chult dotée d'un manche de bois et d'une lame particulièrement affûtée. Elle permettait aussi bien de couper du bois, de se frayer un passage dans la végétation ou de se défendre. Entre des mains expertes, elle pouvait également être lancée.

Le marchand leur proposa enfin un canoë assez grand pour six personnes. Sa fabrication en bois chultien était, selon lui, irréprochable, et les pagaies étaient comprises. Il refusa de le louer, peu convaincu de revoir un jour l'embarcation ou ses occupants, mais accepta d'envisager de le racheter à moitié prix s'ils le rapportaient en bon état.

Après quelques calculs, les aventuriers achetèrent le canoë, deux réserves de répulsif et un récupérateur de pluie. Ils disposaient déjà d'une outre chacun et de potions antipoison. Les provisions ordinaires ne poseraient pas de difficulté immédiate ; l'eau, en revanche, resterait une préoccupation constante pendant le voyage.

Restait à ramener leurs achats jusqu'à l'auberge.

Obélik souleva seul le canoë et prit le chemin de la maison de repos, sous le regard de ses compagnons. Kaya les vit revenir avec une embarcation prévue pour six personnes et leur indiqua les chambres du fond, où ils pourraient la déposer pour la nuit. Le goliath parvint à mener sa charge jusqu'au couloir sans transformer l'auberge en chantier naval.

\begin{quote}\itshape
Sérieusement, les gars... un canoë ?
\end{quote}

Le groupe possédait désormais une guide, un itinéraire et l'équipement nécessaire à son premier départ dans la jungle. Au matin, Azaka viendrait les chercher. Firefinger et son ptéranodon les attendaient quelque part à l'est, avant même qu'ils puissent reprendre la piste du Gardien Bouclier et du mal qui consumait Syndra.

Pour leur première journée de préparatifs, ils avaient au moins accompli l'essentiel.

Ils n'avaient pas oublié le canoë.
